% =====================================================================
%  PART II — ULTRATECH CEMENT LTD
% =====================================================================
\partpage{Part II}{UltraTech Cement Ltd}{The easy-to-value name. Ten years of
audited history, a clean tax line, one dominant product — and a price that
requires a return on capital the company has not earned since FY2021-22.}

\runhead{Part II \textbar\ UltraTech Cement Ltd}

\banner{UltraTech Cement Ltd}{NSE: ULTRACEMCO \textbar\ BSE: 532538}{SELL}{Rs 5,024}{Rs 10,745}{$-$53.2\%}{Standalone core, subsidiaries added}

\begin{paracol}{2}
\section{Investment case and recommendation}

\lead{We initiate on UltraTech Cement with a SELL and an intrinsic value of
Rs 5,024 per share against a traded price of Rs 10,745.} As with Ambuja, five
discounted cash flow methods reconcile exactly. Unlike Ambuja, this is not a company
that destroys value: UltraTech earns above its cost of capital and economic value
added is positive in every forecast year. The question is not whether the business
creates value. It is whether it creates enough to justify 3.3 times invested capital
and 21.1 times EBITDA.

\subsection{The argument in three steps}

\subsubsection{The company buys growth with capital rather than earning it with
price} Volume compounded at 12.5\% a year over ten years while realisation managed
2.0\% and fell 6.1\% in FY2024-25. Revenue growth of 14.7\% a year was substantially
acquired, not organic.

\subsubsection{The acquisitions dilute returns} Management's own disclosure puts core
UltraTech assets at roughly Rs 966 of EBITDA per tonne, against Rs 386 for India
Cements and Rs 755 for Kesoram. Return on invested capital fell from 14.5\% in
FY2021-22 to 10.8\% in FY2025-26 — below our estimated cost of capital of 11.43\%.

\subsubsection{The forecast already assumes that reverses} Our model takes ROIC back
up from 10.8\% to 15.2\% over five years, on management's own guided cost savings and
a utilisation ramp. Even with that reversal fully credited, the value is Rs 5,024.

\begin{keybox}\boxtitle{The discount rate is not the explanation}
Solved backwards, the traded price of Rs 10,745 requires a weighted average cost of
capital of about 9.2\% against our 11.43\%. For a company funded roughly 94\% by
equity, that implies an equity risk premium far below anything defensible for India.
The gap is an operating-assumption gap, not a discount-rate gap.
\end{keybox}

\subsection{Why this is the easy company}

UltraTech publishes a ten-year standalone financial highlights table that ties both
ways in every year: net worth plus loan funds plus leases plus deferred tax equals
net fixed assets plus investments plus liquid investments plus net working capital.
The reported effective tax rate of 26.1\% is materially in line with the statutory
25.168\%, so no migration path is needed. There is one dominant product, disclosed
volume, disclosed capacity, and management guidance on capex and cost savings. Almost
nothing had to be reconstructed.

\switchcolumn
\ex{utcl_football.png}{\textbf{Exhibit U1.} The sell-side operating case is that
house's own driver set — EBITDA per tonne rising 54\% in three years — run through
our discounted cash flow rather than through their exit multiple. Source: authors'
model; peer annual reports; NSE close 30 March 2026.}

\ex{utcl_roic.png}{\textbf{Exhibit U4.} The decisive exhibit. Realised ROIC is
recomputed after tax at the statutory rate on opening invested capital. The forecast
assumes the four-year decline reverses and then some. Source: FY2025-26 integrated
report, ten-year standalone highlights; authors' model.}

\ex{utcl_eva.png}{\textbf{Exhibit U5.} Unlike Ambuja, UltraTech's NOPAT covers its
capital charge throughout. The company creates economic value; the question is
whether it creates enough. Source: authors' model.}

\switchcolumn*
\section{The business and how it is put together}

UltraTech is India's largest cement producer and part of the Aditya Birla group.
Standalone grey cement sales were 140.75 million tonnes on revenue of Rs 82,170 crore
in FY2025-26, at an EBITDA of Rs 15,439 crore. The company holds 74.99\% of the
separately listed India Cements and has absorbed Kesoram's cement assets.

\subsection{Standalone core with subsidiaries added}

The operating forecast is built on the standalone entity because only the standalone
series carries the audited ten-year history, the disclosed grey cement volume of
140.75 million tonnes and the capacity figures on which a capacity-times-utilisation
forecast depends. Consolidated volume and consolidated capital work in progress are
not separately disclosed. The subsidiary block is then carried explicitly: base-year
subsidiary revenue of Rs 6,342 crore at a 24.9\% margin, forecast alongside the core,
with the consolidated asset base and the non-controlling interest both brought into
the model.

\begin{okbox}\boxtitleg{The omission is small, and quantified}
Standalone is 92.8\% of group revenue and 90.7\% of group EBITDA. Consolidated
accounts carry Rs 7,909 crore of goodwill, which is 8.9\% of the asset base against
34.2\% at Ambuja. The contrast with Ambuja is one of degree rather than of kind, and
saying so is more useful than asserting that the standalone basis is simply correct.
\end{okbox}

\begin{warnbox}\boxtitler{A disclosed inconsistency between the two models}
UltraTech holds 74.99\% of the listed India Cements and is modelled from a standalone
core. Ambuja holds 50.05\% of the listed ACC and is modelled consolidated. Two
companies in a similar structural position are treated differently, and the reason is
data availability rather than principle — the irony being that UltraTech's stake is
the more controlling of the two. We state this rather than hide it. See Appendix A
for the full treatment and its quantified effect.
\end{warnbox}

\subsection{An unresolved item in the bridge}

The equity bridge adds investments in subsidiaries and others of Rs 6,740 crore. The
FY2025-26 balance sheet reports Rs 12,541 crore under that heading, of which the India
Cements stake is Rs 8,970 crore at carrying value. The two do not tie. Critically,
India Cements' \emph{operations} are already inside the forecast through the
subsidiary block, so adding the stake at book on top would double-count it — an asset
consistency violation. We have therefore left the bridge as modelled and flagged the
reconciliation as an open item rather than adding Rs 197 per share on an assumption.

\switchcolumn
\ex{utcl_pertonne.png}{\textbf{Exhibit U10.} The acquisition drag, in management's
own numbers. Closing the gap between Rs 386 and Rs 966 is the entire bull case.
Source: management commentary relayed through broker research — flagged for
re-sourcing from the earnings call transcript.}

\begin{extab}{Standalone against consolidated, FY2025-26}{The subsidiary column is
derived as consolidated less standalone. Source: FY2025-26 integrated report,
standalone and consolidated statements.}
\begin{tabular}{@{}p{0.36\linewidth}rrr@{}}
\hrow \wh{Rs crore} & \wh{S/A} & \wh{Consol} & \wh{Subs}\\
Revenue & 82,170 & 88,512 & 6,342\\
\arow EBITDA & 15,439 & 17,020 & 1,581\\
EBITDA margin & 18.8\% & 19.2\% & 24.9\% \\
\arow Borrowings + leases & 20,480 & 23,755 & 3,275\\
Goodwill & — & 7,909 & 7,909\\
\arow Non-controlling interest & — & 4,089 & 4,089\\
\trow S/A \% of group revenue & \multicolumn{3}{r}{\textbf{92.8\%}}\\
\trow S/A \% of group EBITDA & \multicolumn{3}{r}{\textbf{90.7\%}}\\
\end{tabular}
\end{extab}

\ex{utcl_netdebt.png}{\textbf{Exhibit U12.} Net debt had been reduced to Rs 1,994
crore by FY2023-24 before the India Cements and Kesoram acquisitions returned it to
Rs 15,462 crore. Source: FY2025-26 integrated report, ten-year standalone highlights.}

\switchcolumn*
\section{Industry context and competitive position}

UltraTech is the share leader in a regional, freight-constrained commodity industry
with roughly 6--7\% volume demand growth. Its competitive position is genuinely
strong: the largest capacity base, the widest geographic spread, and the scale to
absorb distressed assets when they come to market. None of that is in dispute here.

What is in dispute is the price. The industry has been adding capacity faster than
demand, which is why utilisation across the sector sits near 80\% and why realisation
has been flat in nominal terms and falling in real terms. In that environment, scale
buys resilience rather than pricing power, and acquisitions bought at full prices
dilute returns before they improve them.

\begin{keybox}\boxtitle{The reverse cross-check on the sector multiple}
A 19.05x exit EV/EBITDA — the peer median — implies terminal growth of 9.8\% against
a WACC of 11.43\% and nominal GDP of 10.76\%. That is growth at roughly 91\% of the
whole economy in perpetuity, from a company whose realisation has compounded at 2.0\%
over ten years. The first figure in the reverse table that is arguable is an 8x exit
multiple, which implies 6.8\% growth.
\end{keybox}

\section{Ten years of history: what actually happened}

\subsection{Volume compounded; price did not}

The ten-year standalone series is the strongest evidence in this report, because it
is audited, complete and published by the company itself. Grey cement volume rose
from 48.87 to 140.75 million tonnes, a compound rate of 12.5\%. Realisation per tonne
rose from Rs 4,889 to Rs 5,838, a compound rate of 2.0\%, and fell 6.1\% in
FY2024-25 alone. Revenue therefore compounded at 14.7\% — but that growth was bought.

\subsection{Margin did not follow scale}

EBITDA margin was 20.8\% in FY2016-17 and 18.8\% in FY2025-26. The FY2020-21 peak of
25.4\% was a fuel-cost windfall, not an operating achievement. Ten years of
consolidation, scale and geographic expansion did not move the structural margin.

\subsection{And the return fell}

After-tax return on invested capital, measured on opening capital at the statutory
rate, ran at 14.5\% in FY2021-22 and 10.8\% in FY2025-26. Invested capital rose from
Rs 23,546 crore to Rs 86,090 crore over the decade. The company got much bigger and
somewhat less profitable per rupee employed.

\begin{warnbox}\boxtitler{A live error in the source workbook}
The valuation workbook computes historical NOPAT with a tax reference that resolves
to an empty cell, so its published ten-year ROIC series is a \emph{pre-tax} return —
14.4\% for FY2025-26 rather than the correct after-tax 10.8\%. Exhibit U4 uses the
corrected after-tax series. The correction is listed in Appendix B and is applied in
the accompanying corrected workbook.
\end{warnbox}

\switchcolumn
\ex{utcl_hist.png}{\textbf{Exhibit U6.} Two measures on two panels rather than one
chart with two axes, because the scales are unrelated. Source: FY2025-26 integrated
report, ten-year standalone financial highlights.}

\ex{utcl_revmarg.png}{\textbf{Exhibit U7.} Revenue compounded at 14.7\%; margin ended
two points below where it started. Source: as above.}

\begin{extab}{Ten-year standalone record}{Source: UltraTech Cement Ltd, Integrated
and Sustainability Report 2025-26, Standalone Financial Highlights. FY2018-19 and
FY2022-23 are shown as restated in the company's own table.}
\begin{tabular}{@{}p{0.28\linewidth}rrrr@{}}
\hrow \wh{} & \wh{FY17} & \wh{FY21} & \wh{FY24} & \wh{FY26}\\
Volume (Mn.T) & 48.9 & 81.3 & 112.8 & 140.8\\
\arow Revenue (Rs bn) & 239 & 432 & 686 & 822\\
EBITDA (Rs bn) & 49.7 & 110 & 126 & 154\\
\arow EBITDA margin & 20.8\% & 25.4\% & 18.4\% & 18.8\% \\
Realisation (Rs/t) & 4,889 & 5,315 & 6,085 & 5,838\\
\arow Invested cap. (Rs bn) & 235 & 470 & 653 & 861\\
ROIC after tax & — & 14.0\% & 12.8\% & 10.8\% \\
\arow Net debt (Rs bn) & $-$24 & 47 & 20 & 155\\
Net worth (Rs bn) & 239 & 434 & 591 & 747\\
\end{tabular}
\end{extab}

\switchcolumn*
\section{Forecast assumptions}

The explicit period is five years. UltraTech is already at 82\% utilisation with a
clean, complete history and management guidance on both capex and cost savings; it is
much closer to steady state than Ambuja, and a five-year window reaches it. The cost
is that terminal value carries 78.7\% of enterprise value, which is high, and we say
so.

\subsection{The forecast is built the way the industry is run}

Volume is capacity times utilisation, not a growth rate, because UltraTech has already
built the capacity and is filling roughly 82\% of it — so several years of volume
growth require little new capital. EBITDA is volume times EBITDA per tonne, the metric
management actually steers by. Capex is an explicit input anchored to management's
guidance of roughly Rs 10,000 crore a year.

\begin{okbox}\boxtitleg{The correction that mattered most in the build}
An earlier version derived capex from capital intensity, scaling fixed assets to
revenue. That charged the company a second time for capacity it had already paid for
and understated value by roughly a third. Capex as an explicit input is both more
defensible and considerably more generous to the company.
\end{okbox}

\subsection{Deliberately below the sell side}

Our EBITDA per tonne path rises from Rs 1,097 to Rs 1,379 over five years, capturing
roughly two-thirds of management's guided Rs 300 per tonne cost saving. A sell-side
house models Rs 1,423 by FY2027-28 — a 54\% rise in three years. We take the slower
path, and we test theirs in the scenario table.

\begin{notebox}\boxtitlea{Capex may be light, and the direction is against us}
Total capex of Rs 45,000 crore less maintenance of about Rs 23,700 crore leaves
growth capex of Rs 21,300 crore against capacity addition of 47.5 million tonnes, or
roughly Rs 4,500 per tonne. Indian brownfield cement capacity typically costs
Rs 5,000--7,000 per tonne. Correcting this would lower our value further, so the SELL
is conservative — but the inconsistency is real and is recorded rather than hidden.
\end{notebox}

\switchcolumn
\begin{extab}{Forecast drivers, FY2026-27 to FY2030-31}{Source: authors' assumptions
anchored to management guidance and the FY2025-26 report.}
\begin{tabular}{@{}p{0.46\linewidth}rrr@{}}
\hrow \wh{Driver} & \wh{FY27} & \wh{FY29} & \wh{FY31}\\
Capacity (Mn.T) & 186 & 205 & 220\\
\arow Utilisation & 82.0\% & 84.5\% & 86.0\% \\
Volume (Mn.T) & 152.5 & 173.2 & 189.2\\
\arow EBITDA/t growth & 8.0\% & 4.0\% & 2.5\% \\
EBITDA per tonne (Rs) & 1,185 & 1,306 & 1,379\\
\arow Realisation (Rs/t) & 6,072 & 6,379 & 6,556\\
Subsidiary revenue growth & 6.0\% & 5.0\% & 4.0\% \\
\arow Subsidiary margin & 25.0\% & 25.4\% & 25.8\% \\
Consolidated margin (output) & 19.9\% & 20.8\% & 21.3\% \\
\arow Capex (Rs cr) & 9,500 & 9,000 & 8,500\\
Depreciation, \% opening & 5.05\% & 5.05\% & 5.05\% \\
\arow NWC, \% of revenue & $-$3.74\% & $-$3.74\% & $-$3.74\% \\
\end{tabular}
\end{extab}

\ex{utcl_capacity.png}{\textbf{Exhibit U11.} Peak forecast volume of 189.2 million
tonnes sits inside capacity of 220 million tonnes at all times. Source: authors'
model; announced expansion schedule.}

\ex{utcl_capex.png}{\textbf{Exhibit U14.} Capex runs above depreciation in every
forecast year, so the asset base grows in real terms. Source: authors' model.}

\switchcolumn*
\section{Cost of capital}

The inputs are the same family as Ambuja's: an FBIL 10-year G-Sec par yield of 6.85\%
less the 1.87\% sovereign default spread, giving a default-free rupee rate of 4.98\%;
Damodaran's January 2026 India equity risk premium of 7.08\%; and the cash-corrected
unlevered building-materials beta of 0.90. UltraTech carries meaningfully more debt
than Ambuja — Rs 20,480 crore of borrowings and leases, a market debt-to-equity ratio
of 6.47\% — so the capital structure actually does something here.

\subsection{Relevering under the stated debt policy}

As with Ambuja, the model assumes constant $D/V$ with continuous rebalancing, which
is what permits the tax shield to be discounted at $\rho$. Under that policy the
relever carries no tax term, giving a levered beta of 0.9582 and a cost of equity of
11.76\%. The identity $\rho = \text{WACC} + K_d t (D/V)$ ties exactly.

\subsection{Circularity is disclosed, and it does not bite}

The weights use market capitalisation — the very price this report is testing. That
tension is disclosed rather than hidden, and the iterative alternative was solved. On
the model's own equity value the debt-to-equity ratio rises to roughly 14.8\%, which
pulls WACC down by substituting cheaper debt into the weights while pushing the cost
of equity up through a higher relevered beta. The two effects very nearly cancel and
the converged value differs by about 1\%.

\begin{keybox}\boxtitle{Sensitivity to beta is the honest caveat}
A one-point move in beta changes the cost of equity by roughly 700 basis points and
moves the answer by about a third. Until beta is settled from a full exchange-sourced
price history, the output should be read as a band rather than a point estimate. Our
bottom-up 0.958 sits below the regression estimate of 1.086, which again means the
conservative choice was not the one we took.
\end{keybox}

\switchcolumn
\begin{extab}{Cost of capital, UltraTech}{Source: RBI/FBIL; Damodaran January 2026;
authors' analysis.}
\begin{tabular}{@{}p{0.60\linewidth}r@{}}
\hrow \wh{Input} & \wh{Value}\\
10-year G-Sec par yield (FBIL, w/e 20-Mar-26) & 6.85\% \\
\arow Less: India sovereign default spread & $-$1.87\% \\
\trow \textbf{Risk-free rate} & \textbf{4.98\%} \\
India equity risk premium & 7.08\% \\
\arow Unlevered beta, building materials & 0.900\\
Market debt / equity & 6.47\% \\
\arow Levered beta, $\beta_U(1+D/E)$ & 0.9582\\
\trow \textbf{Cost of equity (CAPM)} & \textbf{11.76\%} \\
Pre-tax cost of debt & 8.34\% \\
\arow Statutory tax rate (s.115BAA) & 25.168\% \\
Debt / enterprise value & 6.08\% \\
\trow \textbf{WACC} & \textbf{11.43\%} \\
\trow \textbf{Unlevered rate $\rho$} & \textbf{11.56\%} \\
\arow Check: $\rho = \text{WACC} + K_d t (D/V)$ & ties\\
Regression beta (60 monthly obs.) & 1.086\\
\end{tabular}
\end{extab}

\begin{extab}{Circularity: market weights against iteration}{Solved per SSRN 413240.
Source: authors' analysis.}
\begin{tabular}{@{}p{0.52\linewidth}rr@{}}
\hrow \wh{} & \wh{Mkt cap} & \wh{Iterated}\\
Debt / equity & 6.47\% & 14.8\% \\
\arow Relevered beta & 0.958 & 1.033\\
Cost of equity & 11.76\% & 12.29\% \\
\arow WACC & 11.43\% & 11.35\% \\
\trow \textbf{Difference in value} & \multicolumn{2}{r}{\textbf{$\sim$1.0\%}}\\
\end{tabular}
\end{extab}

\switchcolumn*
\section{Discounted cash flow valuation}

Free cash flow to the firm is positive in every forecast year, rising from Rs 4,394
crore in FY2026-27 to Rs 14,282 crore in FY2030-31. The first year is depressed by a
Rs 2,150 crore working capital outflow as the negative working capital position
normalises against a larger revenue base.

Terminal growth is 6.99\%, the same 65\% of nominal GDP applied to Ambuja, so the two
companies are treated consistently. Terminal reinvestment equals $g$ times invested
capital, which forces a genuine steady state and is what allows all five methods to
reconcile.

\begin{warnbox}\boxtitler{The terminal return is the weak point of this model}
Terminal ROIC comes out at 15.9\% against a WACC of 11.43\% — a permanent spread of
4.5 percentage points in perpetuity. That assumes supernormal returns forever and
contradicts both competitive convergence and UltraTech's own record of a return that
\emph{fell} from 14.5\% to 10.8\% over four years. Golden Check 6 is recorded as a
\FAIL. Fading the terminal return to WACC plus 0.5\% would take the value from
Rs 5,024 to roughly Rs 4,560. We have not applied the fade in the headline, which
means our valuation is if anything too generous and the SELL is conservative.
\end{warnbox}

\subsection{Result}

\begin{keybox}\boxtitle{Result}
Enterprise value Rs 167,767 crore. Equity value Rs 148,047 crore. Value per share
\textbf{Rs 5,024} against a traded Rs 10,745, a downside of 53.2\%.
\end{keybox}

\switchcolumn
\ex{utcl_fcff.png}{\textbf{Exhibit U3.} FCFF equals NOPAT plus depreciation less
capex less the increase in net working capital. Source: authors' model.}

\ex{utcl_bridge.png}{\textbf{Exhibit U2.} Loan funds, lease liabilities and the
non-controlling interest are all deducted; liquid investments and subsidiary holdings
are added. Source: authors' model; FY2025-26 balance sheet.}

\begin{extab}{DCF summary}{Source: authors' model.}
\begin{tabular}{@{}p{0.58\linewidth}r@{}}
\hrow \wh{Line} & \wh{Rs cr}\\
PV of explicit FCFF, FY27--FY31 & 35,730\\
\arow PV of terminal value & 132,037\\
\trow \textbf{Enterprise value of operations} & \textbf{167,767}\\
Add: liquid investments & 1,384\\
\arow Add: investments in subsidiaries \& others & 6,740\\
Less: loan funds & $-$22,781\\
\arow Less: lease liabilities & $-$974\\
Less: non-controlling interest & $-$4,089\\
\trow \textbf{Equity value} & \textbf{148,047}\\
Shares outstanding (crore) & 29.468\\
\trow \textbf{Value per share (Rs)} & \textbf{5,024}\\
\arow Terminal value \% of EV & 78.7\% \\
Implied terminal EV/EBITDA & 7.52x\\
\end{tabular}
\end{extab}

\ex{utcl_tv.png}{\textbf{Exhibit U13.} Horizon length, not optimism, explains most of
the difference in terminal dominance between the two companies. Source: authors'
models.}

\switchcolumn*
\section{Method equivalence and consistency audit}

All five methods return Rs 5,024 per share, with a maximum difference of zero to six
decimal places. The model was independently reproduced in Python from the input
assumptions, and the EVA identity — opening invested capital plus the present value
of future economic value added — ties to the enterprise value exactly.

\subsection{The five consistencies}

\subsubsection{Assumptions --- \PASS} Growth requires reinvestment throughout. Volume
is capacity times utilisation. Capex is explicit and never held flat while sales grow.
The terminal year is a forced steady state.

\subsubsection{Risk --- \PASS} FCFF at WACC, CCF and APV at $\rho$, FCFE at $K_e$, EVA
at WACC. Harris--Pringle constant rebalancing is stated and is now applied in the
unlever, the relever and the shield alike.

\subsubsection{Asset --- \PART} Other income is excluded from EBITDA so surplus cash
and investments are non-operating and added back once, and the unlevered beta is
cash-corrected for the same reason. \emph{The gap:} the India Cements stake is
carried at book, the disclosed investments total does not tie to the bridge, and the
stake's operations are already inside the subsidiary block. That reconciliation is
open.

\subsubsection{Liability --- \PASS} Loan funds and leases are in the WACC weights and
in the bridge. Deferred tax liabilities of Rs 8,506 crore are in neither the weights
nor invested capital, consistently on both sides. Non-controlling interest is
deducted.

\subsubsection{Multiplier --- \PART} The implied terminal EV/EBITDA of 7.52x against
peers at 19.0--19.7x does not bracket. The reverse cross-check sharpens the question
rather than dissolving it, and we take a documented position.

\switchcolumn
\begin{extab}{Method equivalence}{Source: authors' model.}
\begin{tabular}{@{}p{0.28\linewidth}p{0.20\linewidth}rr@{}}
\hrow \wh{Method} & \wh{Rate} & \wh{Rs/sh} & \wh{Diff}\\
FCFF & WACC & 5,024 & 0\\
\arow CCF & $\rho$ & 5,024 & 0\\
APV & $\rho$ & 5,024 & 0\\
\arow FCFE & $K_e$ & 5,024 & 0\\
EVA & WACC & 5,024 & 0\\
\trow \textbf{Status} & & \multicolumn{2}{r}{\textbf{RECONCILED}}\\
\end{tabular}
\end{extab}

\begin{extab}{Mohanty Appendix 5A: the seven sanity checks}{Source: authors'
assessment.}
\begin{tabular}{@{}p{0.60\linewidth}l@{}}
\hrow \wh{Check} & \wh{Verdict}\\
1. Revenue growth plausible vs GDP & \PASS\\
\arow 2. Capacity supports projected volume & \PASS\\
3. Capex sufficient to build capacity & \PART\\
\arow 4. Margin expansion realism & \PASS\\
5. Working capital sanity & \PASS\\
\arow 6. Terminal ROIC fades toward WACC & \FAIL\\
7. Terminal value dominance & \PASS\\
\end{tabular}
\end{extab}

\begin{warnbox}\boxtitler{Both shortfalls run against the recommendation}
Capex may be light relative to the capacity assumed, and terminal ROIC does not fade.
Correcting either would \emph{lower} the valuation. That is the direction we would
rather be wrong in, and it is why we present the SELL as conservative rather than
aggressive.
\end{warnbox}

\switchcolumn*
\section{Relative valuation}

Against its peers, UltraTech is not obviously mispriced. At 21.1x EV/EBITDA it trades
about 11\% above the peer median of 19.05x; at 40.8x it is below the peer median P/E
of 47.5x. On relative measures the company looks fully but not absurdly valued.

The discounted cash flow says something different: that UltraTech and its peers
together look expensive against the cash they generate. Both statements can be true,
because they answer different questions. Relative valuation asks whether UltraTech is
priced consistently with other Indian cement companies. The DCF asks whether the cash
flows justify the price at all.

\begin{keybox}\boxtitle{The honest conclusion}
UltraTech is reasonably valued relative to Indian cement, and expensive relative to
the cash it generates. Anyone paying Rs 10,745 is underwriting a return on capital
this company has not earned since FY2021-22.
\end{keybox}

\subsection{Why we do not average the two}

Averaging a DCF with a peer multiple would import into our answer the very assumption
we have shown to be arithmetically impossible — that the sector can grow faster than
the economy forever. The multiple is presented as an outcome to be explained, not as
an input to be blended.

\section{Sensitivity, scenarios and catalysts}

The two-way grid opposite is rebuilt from the forecast cash flows. Across a WACC range
of 9.5\% to 13.5\% and terminal growth from 4.0\% to 7.0\%, the value ranges from
Rs 2,925 to Rs 9,281. No cell reaches Rs 10,745.

\subsection{Catalysts}

\emph{Bull side:} completion of the India Cements and Kesoram brand and cost
integration, which management expects to lift those assets from Rs 386 and Rs 755 per
tonne toward the Rs 966 the core already earns; quarterly EBITDA-per-tonne prints
confirming the guided Rs 300 per tonne cost saving; the green power share rising from
42\% toward 65\%; and a cement price recovery after realisation fell 6.1\% in
FY2024-25. \emph{Bear side:} industry capacity additions outrunning demand and holding
utilisation near 80\%, which would keep return on capital pinned at its cost.

The single most informative disclosure is blended EBITDA per tonne. If it reaches
Rs 1,423 by FY2027-28 the gap narrows materially; if it stalls near Rs 1,200 the DCF
view strengthens.

\switchcolumn
\ex{utcl_sens.png}{\textbf{Exhibit U8.} Rebuilt two-way grid, live from the forecast
cash flows. The boxed cell is the base case. The grid in the source workbook was
found to discount depreciation instead of free cash flow and to build its terminal
value off revenue instead of NOPAT; see Appendix B. Source: authors' model.}

\ex{utcl_peers.png}{\textbf{Exhibit U9.} Source: peer annual reports and disclosed
market capitalisations, trailing twelve months to June 2026.}

\begin{extab}{Scenario view}{The sell-side column uses that house's own operating
drivers, run through our DCF rather than their exit multiple. Source: authors' model.}
\begin{tabular}{@{}p{0.34\linewidth}rrr@{}}
\hrow \wh{Driver} & \wh{Bear} & \wh{Base} & \wh{Sell-side}\\
Utilisation by FY31 & 80\% & 86\% & 88\% \\
\arow EBITDA/t by FY31 (Rs) & 1,200 & 1,379 & 1,600\\
Terminal ROIC & 11.9\% & 15.9\% & 17.5\% \\
\arow Capex p.a. (Rs cr) & 10,000 & 9,000 & 9,500\\
\trow \textbf{Value per share (Rs)} & \textbf{$\sim$3,900} & \textbf{5,024} & \textbf{$\sim$6,130}\\
\arow Gap to Rs 10,745 & $-$64\% & $-$53\% & $-$43\% \\
\end{tabular}
\end{extab}

\begin{notebox}\boxtitlea{On the sell-side target}
The Rs 15,000 target we tested against is not a discounted cash flow at all: it is
19 times EV/EBITDA applied to the average of two forecast years, so it assumes the
multiple rather than testing it. The stock fell from Rs 12,370 at the date of that
note to Rs 10,745 by 30 March 2026. We use that house's facts and reject its forecast.
\end{notebox}

\switchcolumn*
\section{Risks and what would change our mind}

\subsubsection{Integration delivers on schedule} If India Cements and Kesoram reach
Rs 900 per tonne within three years rather than gradually, blended unit economics
improve much faster than modelled and the value rises materially.

\subsubsection{A real pricing cycle} Our realisation path grows at half the EBITDA per
tonne growth rate. Cement has seen sustained price recoveries before, and one would
break that assumption in the company's favour.

\subsubsection{Our terminal treatment is too harsh, not too generous} We flag the
absent ROIC fade as a failure. A reader who believes UltraTech genuinely can sustain a
4.5-point spread over its cost of capital forever would regard our headline as
correct and our self-criticism as excessive.

\subsubsection{Beta} At a bottom-up 0.958 against a regression 1.086, we again chose
the figure that raises the value.

\begin{keybox}\boxtitle{What would change our mind}
Two consecutive years of after-tax ROIC above 13\% accompanied by blended EBITDA per
tonne above Rs 1,300 and stable realisation. That would demonstrate that the
acquisitions are being fixed rather than merely absorbed, and would move the value
toward the upper end of our scenario range.
\end{keybox}

\switchcolumn
\begin{extab}{Recommendation summary}{Source: authors' model; NSE close 30 March
2026.}
\begin{tabular}{@{}p{0.58\linewidth}r@{}}
\hrow \wh{Item} & \wh{Value}\\
Recommendation & \textbf{SELL}\\
\arow Intrinsic value per share & Rs 5,024\\
Market price, 30 March 2026 & Rs 10,745\\
\arow Downside & $-$53.2\% \\
Market capitalisation & Rs 316,633 cr\\
\arow Model equity value & Rs 148,047 cr\\
WACC & 11.43\% \\
\arow Terminal growth & 6.99\% \\
Explicit forecast horizon & 5 years\\
\arow Terminal value \% of EV & 78.7\% \\
Implied terminal EV/EBITDA & 7.52x\\
\arow WACC implied by market price & 9.17\% \\
Value with terminal ROIC fade applied & Rs 4,560\\
\end{tabular}
\end{extab}

\begin{okbox}\boxtitleg{Classification: the easy company}
UltraTech is the easy-to-value name: one dominant product, ten years of audited and
internally consistent history, disclosed volume and capacity, an effective tax rate in
line with statutory, and explicit management guidance on capex and cost savings. The
modelling problems it posed were ordinary — choosing a horizon, handling circularity,
deciding the terminal return — rather than structural.
\end{okbox}
\end{paracol}
\clearpage
